\documentclass[11pt,a4paper]{article}
\usepackage[margin=2.3cm]{geometry}
\usepackage{amsmath,amssymb,booktabs,hyperref,enumitem}
\usepackage[T1]{fontenc}
\usepackage{xcolor}
\hypersetup{colorlinks=true,linkcolor=blue!60!black,urlcolor=blue!60!black,citecolor=blue!60!black}
\newcommand{\EbNo}{E_b/N_0}
\newcommand{\EbNJ}{E_b/N_J}
\newcommand{\Pfa}{P_{\mathrm{fa}}}
\newcommand{\Pd}{P_{\mathrm{d}}}
\setlength{\parskip}{4pt}

\title{S3a: Detecting a partial-band jammer from pilot energy}
\author{Chang Young Ham}
\date{September 2026 \quad (accompanies \texttt{src/s3a\_jammer\_detection.py})}

\begin{document}
\maketitle

\noindent\textbf{Purpose.}
S0--S2 treat the jammer as a known input. S3a asks how many hops a receiver
must observe before it can say that a partial-band jammer is present. The
detector is a classical energy test, so every performance number has a closed
form against which the simulator is checked.
Mathematics used: Gaussian likelihoods, the Neyman--Pearson lemma, chi-square
and binomial distributions.

\section{The per-hop statistic}

Each hop carries $n$ known BPSK pilot symbols with amplitude $\sqrt{E_b}=1$,
received as $r_i = s_i + w_i$ exactly as in S0. The receiver subtracts the known
pilots, which leaves noise only. As in S1, a hop lands in the jammed part of the
band with probability $\rho$; there the noise density is $N_0 + N_J/\rho$
instead of $N_0$. With $\sigma^2 = N_0/2$ the per-hop statistic is
\begin{equation}
T = \frac{1}{\sigma^2}\sum_{i=1}^{n}(r_i - s_i)^2,\qquad
\text{unjammed: } T\sim\chi^2_n,\qquad
\text{jammed: } T/a\sim\chi^2_n,
\label{eq:T}
\end{equation}
\begin{equation}
a = 1 + \frac{N_J/N_0}{\rho},\qquad
\frac{N_J}{N_0} = \frac{\EbNo}{\EbNJ},\qquad
\EbNJ\,[\mathrm{dB}] = G_p - J/S\,[\mathrm{dB}],
\label{eq:a}
\end{equation}
where the last two relations are the S1 link budget. The noise ratio $a$ is the
only quantity through which $J/S$ and $\rho$ enter S3a. Energy statistics of
this kind are the classical energy detector: with noise only, the sum of squared
Gaussian samples follows a central chi-square distribution \cite{urkowitz}.
Here both hypotheses are central chi-square with different scales, because the
jammer is modelled as noise.

At the S1 outage point ($J/S = 0.46$\,dB) the noise ratio is $a = 12.7$ for the
jammer's worst-case $\rho = 0.0095$, $a = 2.11$ for $\rho = 0.1$ and
$a = 1.11$ for $\rho = 1$. A jammer that concentrates its power makes each
jammed hop obvious, but few hops are jammed.

\section{Why the energy test is optimal for one hop}

For the $n$ independent residual samples $x_i = r_i - s_i$, the likelihood ratio of ``jammed'' ($N(0,a\sigma^2)$)
to ``unjammed'' ($N(0,\sigma^2)$) is
\begin{equation}
\Lambda(T) = \prod_{i=1}^{n}
\frac{(2\pi a\sigma^2)^{-1/2}\,e^{-x_i^2/(2a\sigma^2)}}{(2\pi\sigma^2)^{-1/2}\,e^{-x_i^2/(2\sigma^2)}}
= a^{-n/2}\exp\!\Big(\frac{T}{2}\big(1-\tfrac{1}{a}\big)\Big).
\label{eq:lr}
\end{equation}
For $a>1$ it increases with $T$, so ``$\Lambda > \eta$'' is the same test as
``$T>\gamma$''. By the Neyman--Pearson lemma this test has the largest detection
probability of all tests with the same false-alarm probability, whatever the
value of $a$. Its performance follows from Eq.~\eqref{eq:T}:
\begin{equation}
\Pfa = P(\chi^2_n > \gamma),\qquad \Pd = P(\chi^2_n > \gamma/a).
\label{eq:roc}
\end{equation}
The receiver is assumed to know $N_0$ (a calibrated noise floor) and the pilots
(perfect channel knowledge, as for coherent BPSK in S0). If the pilots were not
subtracted, the signal energy would enter $T$ and Eq.~\eqref{eq:lr} would no longer
reduce to an energy comparison.

\section{A window of $M$ hops: the count detector}

The receiver flags each hop with $T>\gamma$ and declares ``jammer present'' when
at least $k$ of $M$ hops are flagged. Hops are independent, so the count is
binomial:
\begin{equation}
\text{no jammer: } \mathrm{Bin}(M,\Pfa),\qquad
\text{jammer: } \mathrm{Bin}(M,p_1),\quad p_1 = \rho\,\Pd + (1-\rho)\,\Pfa .
\label{eq:bin}
\end{equation}
For a given $M$ and $\gamma$, $k$ is the smallest count that keeps the window
false-alarm probability at or below $10^{-3}$. $M^\ast$ is the smallest $M$ whose
window detection probability reaches $0.9$, taking the best per-hop threshold
from a grid of 70 values of $\Pfa$ between $10^{-9}$ and $0.5$. Both targets are
assumed values.

The count detector is simple but not optimal when $\rho<1$: the optimal window
test would add up the per-hop mixture log-likelihood ratios
$\sum_j \ln\big(1-\rho+\rho\,\Lambda(T_j)\big)$, which requires knowing $\rho$ and
$a$. $M^\ast$ is therefore what \emph{this} detector needs; an optimal detector
needs at most as many hops.

\section{Two limits used as checks}

\textbf{Strong jammer (verification 3).} If every jammed hop is flagged
($\Pd\to1$) and false alarms are negligible ($\Pfa\to0$, so $k=1$), the window
only has to contain one jammed hop:
\begin{equation}
1-(1-\rho)^{M} \ge 0.9 \quad\Longleftrightarrow\quad
M \ge \frac{\ln 0.1}{\ln(1-\rho)} \approx \frac{2.30}{\rho}.
\label{eq:limit}
\end{equation}
At $J/S = 30$\,dB the computed $M^\ast$ must be within one hop of the ceiling of
this bound. It can be one hop \emph{below} it because the allowed false alarms
($\Pfa \lesssim 10^{-3}/M$ per hop) add slightly to $p_1$.

\textbf{Full-band jammer (verification 4).} At $\rho=1$ every hop is jammed and
all $Mn$ residual samples are i.i.d.; the likelihood ratio of the whole window has
the form of Eq.~\eqref{eq:lr} with the pooled energy
$\sum_j T_j\sim\chi^2_{Mn}$ (scaled by $a$ when jammed). The pooled-energy test is
therefore optimal there, and by the Neyman--Pearson lemma the count detector,
whose false-alarm probability is at most the same $10^{-3}$, cannot have a higher
detection probability at any $M$. A violation would mean an error in
Eq.~\eqref{eq:roc} or \eqref{eq:bin}.

\section{Monte Carlo checks and the $z$-score}

Verification 1 generates pilots, adds noise, subtracts the pilots and compares
the observed $\Pfa$ and $\Pd$ with Eq.~\eqref{eq:roc}. Verification 2 simulates
whole windows and compares $P(\text{count}\ge k)$ with Eq.~\eqref{eq:bin}, with
and without the jammer. Each point is tested by
\begin{equation}
z = \frac{K - Np}{\sqrt{Np(1-p)}},
\end{equation}
where $K$ is the number of events in $N$ trials. S0 used $\sqrt{Np}$ in the
denominator, the small-$p$ form; here $p$ can be close to 1, so the full binomial
variance is needed. The pass criterion is $|z|<4$ at every point, as in S0 and S1.

\section{What the numbers show}

With $n=16$ pilots (\texttt{results/s3a\_summary.txt}):
\begin{itemize}[leftmargin=*,itemsep=2pt]
\item At the S1 outage point and the jammer's worst-case $\rho = 0.0095$, per-hop
detection is almost certain ($a=12.7$), so $M^\ast = 242$ equals the
strong-jammer bound~\eqref{eq:limit}: the delay is set by how rarely a hop lands
in the jammed band. More pilots do not help ($M^\ast = 242$ for $n=16$ and $32$,
261 for $n=8$). If the jammer must occupy at least one channel ($\rho = 0.01$),
$M^\ast = 230$.
\item At the same $J/S$ with $\rho=1$, each hop is only weakly affected
($a=1.11$). The count detector needs 337 hops and the optimal pooled-energy test
212; here more pilots help roughly in proportion (671, 337, 170 hops for
$n = 8, 16, 32$).
\item Between the two, the count detector needs the fewest hops (92 on the grid,
near $\rho\approx0.04$). This holds for this detector only (Section~3).
\item At $J/S = 12$\,dB, the S1 worst-case $\rho = 0.135$ needs 16 hops, again the
bound~\eqref{eq:limit}; at $\rho=1$ it needs 4 (pooled energy 3).
\end{itemize}
All $M^\ast$ are in hops; converting them to seconds needs a hop rate, which
is not modelled here.

\section{Correspondence with the code}

\begin{center}
\begin{tabular}{@{}ll@{}}
\toprule
Mathematics & Code \\
\midrule
$a$, Eq.~\eqref{eq:a} & \texttt{noise\_ratio()} \\
$T$, Eq.~\eqref{eq:T} & \texttt{simulate\_hop\_energy()} \\
$\gamma$, $\Pd$ from $\Pfa$, Eq.~\eqref{eq:roc} & \texttt{hop\_detection()} \\
smallest admissible $k$ & \texttt{count\_threshold()} \\
$M^\ast$ over the threshold grid & \texttt{hops\_to\_detect()} \\
pooled energy at $\rho=1$ & \texttt{pooled\_detection()}, \texttt{hops\_to\_detect\_pooled()} \\
Eq.~\eqref{eq:limit} & \texttt{strong\_jammer\_limit()} \\
S1 outage point and worst-case $\rho$ & imported from \texttt{s1\_fh\_partial\_band\_jamming.py} \\
\bottomrule
\end{tabular}
\end{center}

\begin{thebibliography}{9}
\bibitem{urkowitz} H.~Urkowitz, ``Energy detection of unknown deterministic
  signals,'' \emph{Proceedings of the IEEE}, vol.~55, no.~4, pp.~523--531, Apr.\ 1967.
\end{thebibliography}

\end{document}
